\chapter{按失效条件分类的反例索引}
\label{app:counterexample-index}
\label{chap:appendix-counterexamples}

\begin{chapterguide}
反例的作用是定位命题中不能删除的条件。本附录按“失效条件”而非按章节排列实例，并区分已由前三编完整验证的条目与依赖后续正文的预留条目。索引中的每个反例都应同时核对保留条件、删去条件和失败结论。
\end{chapterguide}

\section{有界但不收敛、Cauchy 但无极限}

\begin{description}[style=nextline,leftmargin=2.8em]
 \item[有界但不收敛]
 数列 \(a_n=(-1)^n\) 有界，却有极限分别为 \(1,-1\) 的奇偶子列。它表明有界性只能保证存在收敛子列，不能保证全列收敛。
 \item[单调但不收敛于有限数]
 \(a_n=n\) 单调递增却无上界，故趋于 \(+\infty\)。单调有界收敛定理中的有界条件不能省略。
 \item[相邻差趋零但非 Cauchy]
 调和部分和 \(H_n=\sum_{k=1}^n1/k\) 满足 \(H_{n+1}-H_n\to0\)，而
 \[
  H_{2n}-H_n\ge\frac12.
 \]
 Cauchy 条件必须统一控制任意两个充分靠后的项。
 \item[Cauchy 但值域中无极限]
 \chapterref{chap:real-number-system}逼近 \(\sqrt2\) 的有理数列在 \(\Q\) 中是 Cauchy 列，却没有有理极限。失败条件是值域不完备；同一列在 \(\R\) 中收敛。
 \item[无界但有收敛子列]
 定义 \(a_{2n}=0\)、\(a_{2n-1}=n\)。原列无界，但偶数子列收敛到 \(0\)。Bolzano--Weierstrass 定理只给有界列的充分条件，不给“有收敛子列则原列有界”的逆命题。
\end{description}

\section{逐点收敛但非一致收敛}

\begin{remark}[预留说明]
函数列的逐点收敛、一致收敛及交换极限定理安排在第五编“函数列与一致收敛”。该编完成前，本节保留分类位置，不在此先行引入尚未建立的定义。
\end{remark}

后续应收录的标准类型包括：极限在移动尖峰处失去统一控制、连续函数列逐点极限不连续、逐点收敛不能保持积分或微分，以及紧集上一致控制恢复结论的正面边界。正式条目须在相应主章节定稿并建立稳定标签后补入引用。

\section{连续、可导、光滑和解析之间的反例}

前三编已经建立的连续性边界包括：
\begin{description}[style=nextline,leftmargin=2.8em]
 \item[连续但非一致连续]
 \(x^2\) 在 \(\R\) 上连续，点对 \(n,n+1/n\) 的距离趋零而像距趋于 \(2\)。
 \item[一致连续但非 Lipschitz]
 \(\sqrt x\) 在 \([0,\infty)\) 上满足
 \[
  \abs{\sqrt x-\sqrt y}\le\sqrt{\abs{x-y}},
 \]
 因而一致连续；若为 Lipschitz，则 \(x^{-1/2}\) 型差商会在 \(0\) 附近被统一常数控制，矛盾。
 \item[单调但不连续]
 取整函数在每个整数处跳跃。单调性保证左右极限存在，却不保证两侧极限相等。
 \item[严格单调但像集有缺口]
 在 \(\R\) 上定义 \(f(x)=x\)（\(x<0\)）、\(f(x)=x+1\)（\(x\ge0\)）。它严格递增而在 \(0\) 跳跃，像集漏掉 \((0,1)\)。
\end{description}

\begin{remark}[后续预留]
“可导但导函数不连续”“处处可导但非 \(C^1\)”“光滑但非解析”“各阶导数在一点全为零但函数非零”等条目依赖第四编微分学和第七编复分析。相关定义、证明与章节交叉引用暂留待补。
\end{remark}

\section{各种极限与积分、微分不能交换的反例}

当前可完整核对的是代数极限的未定式：
\begin{itemize}
 \item \(f,g\to+\infty\) 时，\(f-g\) 可趋于有限数、\(+\infty\)、\(-\infty\) 或无极限；
 \item \(f\to0\)、\(g\to+\infty\) 时，\(fg\) 可趋于 \(0\)、非零常数、无穷或无极限；
 \item 复合删心极限若内层函数持续取外层中心，外层点值可破坏复合结论。
\end{itemize}

\begin{remark}[后续预留]
极限与微分的交换依赖第四、五编；极限与 Riemann、Lebesgue 积分的交换依赖第五、六编。正式索引将按“缺少一致收敛”“缺少支配”“缺少一致可积性”“缺少导数一致控制”分栏补入。
\end{remark}

\section{紧致性、完备性和闭有界性的反例}

\begin{description}[style=nextline,leftmargin=2.8em]
 \item[有界但不闭]
 \((0,1)\) 的覆盖 \(\{(1/n,1):n\ge2\}\) 没有有限子覆盖；连续函数 \(x\) 在其上不取得最大、最小值。
 \item[闭但无界]
 \(\R\) 的覆盖 \(\{(-n,n)\}\) 没有有限子覆盖；连续函数 \(x^2\) 不一致连续。
 \item[完备子空间必须闭]
 \(\Q\) 作为 \(\R\) 的子空间不闭，逼近 \(\sqrt2\) 的有理 Cauchy 列无有理极限。
 \item[实线中的特殊等价]
 在 \(\R\) 中，闭且有界等价于紧致。该结论不能直接推广为“任意度量空间中闭且有界即紧致”；无限维和一般度量空间的反例留待相应章节。
 \item[连续像的方向]
 连续像保持紧致与连通，但原像一般不保持这些性质；只有闭集原像等局部拓扑结论可直接由连续性得到。
\end{description}

\section{各种测度收敛和 \texorpdfstring{\(L^p\)}{Lp} 收敛的反例}

\begin{remark}[预留：依赖第六编]
本节将在测度、几乎处处收敛、依测度收敛、\(L^p\) 收敛和弱收敛建立后补写。预定分类为：蕴含方向失败、缺少有限测度条件、缺少一致可积性、子列结论不可逆以及端点指数失败。
\end{remark}

\section{有限维结论在无限维空间中的失败}

\begin{remark}[预留：依赖第十四、十五编]
预定条目包括闭有界集非紧、单位球不序列紧、线性算子有界性与矩阵范数类比的边界、谱可能含连续部分，以及有限维范数等价不能无条件统一到维数变化的空间族。
\end{remark}

\section{一般拓扑中序列刻画失效的反例}

\begin{remark}[预留：依赖第八编]
前三编在实直线和度量空间中大量使用序列刻画。一般拓扑空间未必第一可数，闭包、连续性与紧致性可能不能只用序列描述。本节将在网、滤子和可数性公理完成后补入正式反例。
\end{remark}

\section{复分析、概率与算子理论中的边界例}

\begin{remark}[预留：依赖第七、十一、十五编]
本节预定按解析延拓、逐点概率收敛、几乎必然收敛、无界算子定义域和谱映射等主题分类。为避免使用尚未固定的符号与版本，当前仅保留位置。
\end{remark}

\section{弱解、正则性和非线性极限中的反例}

\begin{remark}[预留：依赖第十七至二十编]
预定条目包括弱极限不能保持非线性乘积、弱解不自动光滑、能量界不足以给强收敛、边界正则性失败和尺度临界现象。正式内容须与 Sobolev 空间、分布和具体方程章节同步。
\end{remark}

\section*{使用方法}
\addcontentsline{toc}{section}{使用方法}

检验一个命题时，可按以下顺序使用本索引：
\begin{enumerate}
 \item 把每个假设改写为可逐项核对的定义；
 \item 一次只删除一个条件，并保留其余条件；
 \item 核对反例确实使结论失败，而非仅使原证明失效；
 \item 记录失败机制属于无界、缺边界、不完备、非一致、非紧致或非连通中的哪一类；
 \item 若条目标为预留，不把它当作正文已证明结果引用。
\end{enumerate}

\section*{习题}
\addcontentsline{toc}{section}{习题}

\begin{enumerate}[label=\textbf{D.\arabic*.}]
 \exerciseitem{D.1} 为“有界但不收敛”给出两个机制不同的数列，并分别指出其收敛子列。
 \exerciseitem{D.2} 构造相邻差趋零但非 Cauchy 的数列，并用双指标差证明。
 \exerciseitem{D.3} 解释有理 Cauchy 列无有理极限时，失败的是哪一条假设而不是 Cauchy 定义。
 \exerciseitem{D.4} 构造无界但有两个不同有限部分极限的数列。
 \exerciseitem{D.5} 为连续但不一致连续给出“逃向无穷远”和“逼近缺失边界”两类反例点对。
 \exerciseitem{D.6} 证明 \(\sqrt x\) 一致连续但非 Lipschitz。
 \exerciseitem{D.7} 给出严格单调但不连续的函数，并求其像集。
 \exerciseitem{D.8} 为 \(\infty-\infty\) 与 \(0\cdot\infty\) 各构造至少三种结果。
 \exerciseitem{D.9} 分别说明 \((0,1)\) 与 \(\R\) 的非紧致机制。
 \exerciseitem{D.10} 证明 \(\Q\) 在通常距离下不完备，并指出它在 \(\R\) 中不闭。
 \projectexerciseitem{D.11} \ 【前置：前三编；预计用时：75分钟】选取一个主定理，制作“删条件—反例—失败步骤—可恢复结论”四列表。
 \openexerciseitem{D.12} 说明为什么预留条目在依赖章节完成前不宜写成正式反例索引；至少讨论符号、定理版本和逻辑循环三类风险。
\end{enumerate}
